% problem-000355 1 核反应与元素
\section*{1. 核反应与元素}
\textit{下列为分问。}

\subsection*{1-1}
2009年10⽉合成了第117号元素，从此填满了周期表第七周期所有空格，是元素周期系发展的⼀个⾥程碑。117号元素是⽤249Bk轰击 $^ { 4 8 } \mathrm { C a }$ 靶合成的，总共得到6个117号元素的原⼦，其中1个原⼦经�次α衰变得到270Db后发⽣裂变；5个原⼦则经�次α衰变得到 $^ { 2 8 1 } \mathrm { R g }$ 后发⽣裂变。⽤元素周期表上的117号元素符号，写出得到117号元素的核反应⽅程式（在元素符号的左上⻆和左下⻆分别标出质量数和原⼦序数）。

\subsection*{1-2}
写出下列结构的中⼼原⼦的杂化轨道类型：

（图：）

（图：）

$
[ (\mathrm{C} _ {6} \mathrm{H} _ {5}) \mathrm{IF} _ {5} ] ^ {-}
$

（图：）

$
(\mathrm{C} _ {6} \mathrm{H} _ {5}) _ {2} \mathrm{Xe}
$

$
\mathrm{I} (\mathrm{C} _ {6} \mathrm{H} _ {5}) _ {2} ] ^ {+}
$

\subsection*{1-3}
⼀氯⼀溴⼆(氨基⼄酸根)合钴(III)酸根离⼦有多种异构体，其中之⼀可⽤如下简式表⽰。请依样画出其他所有⼋⾯体构型的异构体。

（图：）
